% Generated draft: requires mathematical and source review.
\documentclass[11pt]{article}
\usepackage[margin=25mm]{geometry}
\usepackage{fontspec}
\setmainfont{Latin Modern Roman}
\usepackage{amsmath,amssymb,mathtools,unicode-math}
\setmathfont{Latin Modern Math}
\usepackage{enumitem,needspace,xurl,hyperref}
\setlength{\parindent}{0pt}
\setlength{\parskip}{5pt}
\setlength{\emergencystretch}{3em}
\newcommand{\N}{\mathbb{N}}
\newcommand{\Z}{\mathbb{Z}}
\newcommand{\Q}{\mathbb{Q}}
\newcommand{\R}{\mathbb{R}}
\newcommand{\C}{\mathbb{C}}
\newcommand{\F}{\mathbb{F}}
\newcommand{\A}{\mathbb{A}}
\newcommand{\PP}{\mathbb{P}}
\newcommand{\E}{\mathbb{E}}
\newcommand{\eps}{\varepsilon}
\newcommand{\dd}{\,\mathrm{d}}
\begin{document}
\section*{UnsolvedMath ID: 20001966}
\textbf{Source problem code:} catalogue entry
\subsection{Definitions and mathematical statement}
Hausel's Conjecture 6.14 (2000) compares a hyperkähler quotient by \(G\) with one by a maximal torus \(T\). Write \(X_G=T^*A/\!/\!/\!/G\), \(X_T=T^*A/\!/\!/\!/T\), \(W=N_G(T)/T\), and let \(\Delta\) be the full root system. Subject to both quotients being ``hyper-compact'', the source proposes
\[
 H^*(X_G)\cong H^*(X_T)^W/\operatorname{Ann}(\widetilde e),
 \qquad \widetilde e=\prod_{\alpha\in\Delta}\alpha
 \in(\operatorname{Sym}\mathfrak t^*)^W\cong H_T^*(\mathrm{pt})^W.
\]
In the quotient formula, \(\widetilde e\) must mean the image of this class in the invariant cohomology ring, and its annihilator consists of invariant classes whose product with that image is zero. The excerpt does not give the coefficient ring, level, regularity hypotheses, or definition of hyper-compactness. In particular, the symmetric-algebra identification requires a coefficient convention; the supplied partial result below is over \(\Q\).

The aid's special case uses \(G=U(k)\), \(A=\operatorname{Hom}(\C^k,\C^n)\), with \(1\le k\le n\), and a positive equal central regular level. It identifies
\[
 X_G\cong T^*\operatorname{Gr}(k,n),\qquad
 X_T\cong(T^*\PP^{n-1})^k,\qquad W\cong S_k.
\]
Set
\[
 R=\Q[h_1,\ldots,h_k]/(h_1^n,\ldots,h_k^n),\quad
 S=R^{S_k},\quad \deg h_i=2,\quad
 \delta=\prod_{1\le i<j\le k}(h_i-h_j).
\]
Here \(S_k\) permutes the \(h_i\). The all-root class is
\[
 \prod_{i\ne j}(h_i-h_j)=(-1)^{k(k-1)/2}\delta^2.
\]
Although \(\delta\) is alternating, \(\operatorname{Ann}_S(\delta)=\{s\in S:s\delta=0\text{ in }R\}\) is well-defined. The reported algebraic claim is
\[
 \operatorname{Ann}_S(\delta)=\operatorname{Ann}_S(\delta^2),\qquad
 S/\operatorname{Ann}_S(\delta^2)
 \cong\Q[c_1,\ldots,c_k]/(b_{n-k+1},\ldots,b_n)
 \cong H^*(\operatorname{Gr}(k,n);\Q).
\]
To fix a sign convention left implicit by the aid, take \(c_i\) to be the elementary symmetric classes and define \(b_j\) by
\[
 (1+c_1t+\cdots+c_kt^k)^{-1}=\sum_{j\ge0}b_jt^j.
\]
Changing \(b_j\) by the usual complete-symmetric sign convention does not change the indicated ideal. More precisely, the ideal statement concerns the kernel of the composite map from \(\Q[c_1,\ldots,c_k]\) to \(S/\operatorname{Ann}_S(\delta^2)\); it is not an identification of an ideal in \(S\) with a polynomial ideal without taking preimages. The aid attributes a direct proof for \(\delta^2\) to a complementary-partition top-coefficient pairing.

This verifies the reported ordinary all-root ring identity for this family only under the aid's later semiprojective/Morse--Bott interpretation of hyper-compactness. It expressly warns that the AIM wording uses finitely many critical points. Cotangent fiber scaling fixes the entire zero section, so finite critical components and finite critical points must not be conflated. The excerpt does not reproduce the full historical definition; no verification under that literal stronger wording is asserted.

The extracted record additionally contains the next Section 6.4 and Theorem 6.15, attributed to Hausel--Szenes (2003), motivated by Witten's work on other moduli spaces. This is separate from Conjecture 6.14. It introduces \(M=M^1_{\mathrm{Dol}}(SL(2,\C))\), a positive integral generator \(\alpha\in H^2_{U(1)}(M)\cong\Z\), and an expression \(BA(y)\), then states a residue pattern
\[
 \int_M e^\alpha=
 \operatorname{Res}_{y=0}BA(y)+\operatorname{Res}_{y=-u}BA(y)
 +\operatorname{Res}_{y=u}BA(y)
 =-\sum_b\operatorname{Res}_{y=b}BA(y),
\]
where \(b\) ranges over solutions of the displayed Bethe--Ansatz equations. The extraction corrupts the fractions and exponent placement in \(BA\) and those equations; the curve, genus assumptions, equivariant parameter \(u\), integration convention, and interpretation of that sum are not supplied. This neighboring formula is preserved only as reported context, not repaired into a new exact theorem or used as a hypothesis of the all-root conjecture.
\subsection{Short English statement}
The conjecture predicts a nonabelian hyperkähler quotient's cohomology from Weyl-invariant torus-quotient cohomology by annihilating the full root product. The aid verifies a rational cotangent-Grassmannian family, while flagging that its meaning of hyper-compactness differs from the historical wording. A damaged residue theorem appended by extraction is separate context.
\subsection{Sources}
\begin{itemize}
\item Local UnsolvedMath registry, permanent ID 20001966, exact source code catalogue entry; canonical \texttt{aim-geometry-notes.json}, \texttt{notes[303]}, tag \texttt{conjecture}; AIM Workshop Problem Lists, Geometry. Created and updated 14 August 2026.
\item AIM \emph{Moment Maps and Surjectivity in Various Geometries}, Conjecture 6.14 (Hausel 2000); \url{https://aimath.org/WWN/momentmaps/momentmaps.pdf}. The record also extracts Section 6.4, Theorem 6.15 (Hausel--Szenes 2003), and mentions motivation from Witten. No full bibliographic entries for these attributions are supplied. The PDF was not accessed.
\item Supplied attempt 1: \texttt{partial\_result}, problem status at run \texttt{partially\_solved}, rigor label \texttt{proved}, novelty confidence \texttt{low}; classification \texttt{PARTIAL-PROGRESS}, suggested difficulty L4, registry status \texttt{partially\_solved}. Candidate: rational Vandermonde/all-root annihilator comparison and cotangent-Grassmannian identity. These are source assessments requiring review.
\item Ulam AI / UnsolvedMath Contributors; CC BY 4.0. Machine-generated research aid, not peer reviewed; independently verify before citation.
\end{itemize}
\end{document}
